% Zone „Das kennst du schon“ – aus bank/orthogonalitaet/zone.jsonl (zusammenbau v0.1)
\einheitenkopf[zone][Kennst du schon]{Das kennst du schon}
\setcounter{aufgabe}{0}

%% TODO Titel: Ich-kann-Satz fehlt in der Bank (Platzhalter: Kettenname) – Nr. 1
%% TODO Anweisung: ein Satz über der ersten Teilaufgabe fehlt in der Bank – Nr. 1
\begin{aufgabe}{Skalarprodukt berechnen und „null heißt senkrecht“ anwenden}
\begin{teile}
\teil Berechne $(2 | 1 | 3) \circ (1 | 4 | -2)$. Stehen die Vektoren senkrecht? \leerfeld
\teil Berechne $(0{,}5 | 2 | -1) \circ (4 | -1{,}5 | 1)$. \leerfeld
\end{teile}
\end{aufgabe}

%% TODO Titel: Ich-kann-Satz fehlt in der Bank (Platzhalter: Kettenname) – Nr. 2
%% TODO Anweisung: ein Satz über der ersten Teilaufgabe fehlt in der Bank – Nr. 2
\begin{aufgabe}{Richtungsvektoren bilden und Beträge berechnen}
\begin{teile}
\teil $A(1 | 2 | 0)$, $B(4 | 6 | 0)$. Gib $\overrightarrow{AB}$ an und berechne $|\overrightarrow{AB}|$. \leerfeld
\teil $C(-1 | 2 | 5)$, $D(3 | -2 | 7)$. Gib $\overrightarrow{CD}$ an und berechne $|\overrightarrow{CD}|$. \leerfeld
\end{teile}
\end{aufgabe}

%% TODO Titel: Ich-kann-Satz fehlt in der Bank (Platzhalter: Kettenname) – Nr. 3
%% TODO Anweisung: ein Satz über der ersten Teilaufgabe fehlt in der Bank – Nr. 3
\begin{aufgabe}{Lineare und quadratische Gleichungen sowie kleine Gleichungssysteme lösen, Lösungen gegen einen Bereich sieben}
\begin{gleichungsraster}[2]
\gl{\text{Löse: $3t - 12 = 0$}} &
\gl{\text{Löse das Gleichungssystem: $2x - y = 4$ und $x + 2y = 7$}} \\
\end{gleichungsraster}
\end{aufgabe}

%% TODO Titel: Ich-kann-Satz fehlt in der Bank (Platzhalter: Kettenname) – Nr. 4
%% TODO Anweisung: ein Satz über der ersten Teilaufgabe fehlt in der Bank – Nr. 4
\begin{aufgabe}{Kollinearität prüfen (Vielfaches)}
\begin{teile}
\teil Sind $(2 | 6 | -4)$ und $(1 | 3 | -2)$ kollinear? Gib den Faktor an. \leerfeld
\teil Sind $(1{,}5 | -3 | 4{,}5)$ und $(1 | -2 | 3)$ kollinear? Gib den Faktor an. \leerfeld
\end{teile}
\end{aufgabe}

%% TODO Titel: Ich-kann-Satz fehlt in der Bank (Platzhalter: Kettenname) – Nr. 5
%% TODO Anweisung: ein Satz über der ersten Teilaufgabe fehlt in der Bank – Nr. 5
\begin{aufgabe}{Normalenvektor aus der Koordinatengleichung ablesen, Spannvektoren aus der Parameterform lesen}
\begin{teile}
\teil $E\colon 2x - 3y + z = 5$. Gib einen Normalenvektor von $E$ an. \leerfeld
\teil $E\colon -x + 2y - 5z + 7 = 0$. Gib einen Normalenvektor von $E$ an. \leerfeld
\end{teile}
\end{aufgabe}

%% TODO Titel: Ich-kann-Satz fehlt in der Bank (Platzhalter: Kettenname) – Nr. 6
%% TODO Anweisung: ein Satz über der ersten Teilaufgabe fehlt in der Bank – Nr. 6
\begin{aufgabe}{Ähnlichkeit rechtwinkliger Dreiecke und Seitenverhältnisse}
\begin{teile}
\teil Zwei rechtwinklige Dreiecke sind ähnlich. Im ersten sind die Katheten 3 cm und 4 cm lang, im zweiten ist die kürzere Kathete 6 cm lang. Wie lang ist die längere Kathete des zweiten Dreiecks? \leerfeld[cm]
\teil Im rechtwinkligen Dreieck teilt die Höhe auf die Hypotenuse diese in Abschnitte von 4 cm und 9 cm. Wie lang ist die Höhe? (Die beiden Teildreiecke sind ähnlich.) \leerfeld[cm]
\end{teile}
\end{aufgabe}

%% TODO Titel: Ich-kann-Satz fehlt in der Bank (Platzhalter: Kettenname) – Nr. 7
%% TODO Anweisung: ein Satz über der ersten Teilaufgabe fehlt in der Bank – Nr. 7
\begin{aufgabe}{Richtungsvektoren bilden und Beträge berechnen – Fehler finden}
\begin{teile}
\teil Für $A(2 | 3 | 1)$ und $B(5 | 1 | 4)$ schreibt Lea: \rechnung{\overrightarrow{AB} &= (2 - 5 \mid 3 - 1 \mid 1 - 4) = (-3 \mid 2 \mid -3)} Finde den Fehler.
\end{teile}
\end{aufgabe}

%% TODO Titel: Ich-kann-Satz fehlt in der Bank (Platzhalter: Kettenname) – Nr. 8
%% TODO Anweisung: ein Satz über der ersten Teilaufgabe fehlt in der Bank – Nr. 8
\begin{aufgabe}{Richtungsvektoren bilden und Beträge berechnen – selbst rechnen}
\begin{teile}
\teil $A(-1 | 4 | 2)$, $B(3 | 1 | 2)$. Gib $\overrightarrow{AB}$ an und berechne seine Länge. \leerfeld
\end{teile}
\end{aufgabe}
